\begin{french}

\chapter*{Résumé étendu en français}
\addcontentsline{toc}{chapter}{Résumé étendu en français}
\chaptermark{Résumé français}
\markboth{Résumé français}{Résumé français}

\emph{Ce résumé est une version abrégée du \Cref{ch:intro}. Nous recommandons plutôt la lecture dudit chapitre pour avoir une compréhension plus globale de la thèse.}

\section*{Les deux facettes du problème du morphisme}

Cette thèse s'intéresse à des variations du \emph{problème du morphisme}, qui est un concept 
central en informatique théorique. La notion de \emph{morphisme} formalise l'idée de fonction
préservant de l'information relationnelle d'une structure à une autre. Ces objets, ainsi que le
problème de décision associé, apparaissent naturellement dans divers domaines : citons par exemple
en \emph{logique}, en tant que caractérisation du \emph{model-checking} pour les formules 
primitives positives ; en \emph{théorie des bases de données}, comme la sémantique des requêtes 
conjonctives ; et en \emph{programmation logique}, comme formalisation du fait qu'une
structure satisfait un ensemble de contraintes.

Dans cette section, nous commençons par formaliser la notion de morphisme, avant
de donner deux exemples représentatifs de l'utilisation des problèmes de morphismes en informatique
fondamentale : l'évaluation de requêtes SQL et la résolution de Sudoku.
Ces deux exemples révèlent une \emph{dualité} sur la façon dont les problèmes sont encodés : dans
certains cas, la requête---ou formule logique---apparaît à gauche du morphisme, alors que dans d'autres cas elle apparaît à droite.
Le premier type d'encodage correspond à ce que nous appelons des \emph{problèmes existentiels},
qui forment les fondements de la \emph{théorie des bases de données}, que nous présenteront dans la seconde section. Quant au second type, il encode des \emph{problèmes universels}, étudiés
dans le domaine des \emph{problèmes de satisfaction de contraintes}, présenté dans la dernière partie de ce résumé. Cette dualité intrinsèque du problème du morphisme---ou plutôt de l'encodage
de problèmes naturels dans le problème du morphisme---guidera cette thèse jusque dans sa structure, que l'on peut notamment observer par sa division en deux parties.

\begin{figure}
	\centering
	\begin{tikzpicture}
		\foreach \i in {0,1,...,5}{
			\node[vertex] (\i) at (\i*360/6: 1.2cm) {};
		}

		\draw[edge] (1) to (0);
		\draw[edge] (2) to (1);
		\draw[edge] (3) to (2);
		\draw[edge] (3) to (4);
		\draw[edge] (5) to (4);
		\draw[edge] (5) to (0);

		\node[vertex] at (4.33,.66) (a) {};
		\node[vertex] at (5.66,.66) (b) {};
		\node[vertex] at (5.66,-.66) (c) {};
		\node[vertex] at (4.33,-.66) (d) {};
		\node[vertex] at (6.8,0) (e) {};

		\draw[edge] (d) to (a);
		\draw[edge] (a) to (b);
		\draw[edge] (b) to (c);
		\draw[edge] (d) to (c);
		\draw[edge] (b) to (e);
		\draw[edge] (e) to (c);

		\draw[edge, draw=cBlue, dotted, out=-20, in=160] (3) to (d);
		\draw[edge, draw=cBlue, dotted, out=-25, in=170] (2) to (a);
		\draw[edge, draw=cBlue, dotted, out=5, in=130] (1) to (b);
		\draw[edge, draw=cBlue, dotted, out=10, in=145] (0) to (c);
		\draw[edge, draw=cBlue, dotted, out=-20, in=195] (5) to (d);
		\draw[edge, draw=cBlue, dotted, out=-30, in=210] (4) to (c);
	\end{tikzpicture}
	\caption{
		\AP\label{fig:example-graph-homomorphism-fr}
		Deux graphes (en noir) et un homomorphisme (flèches pointillées bleues) du graphe de gauche vers celui de droite.
		(Reproduction de \Cref{fig:example-graph-homomorphism}.)
	}
\end{figure}
La structure mathématique la plus simple est peut-être celle d’un graphe fini (orienté) :
elle consiste en un ensemble fini $V$ de sommets (également appelé \emph{domaine}),
ainsi qu’en un ensemble d’\emph{arêtes} $\+E \subseteq V \times V$.
Un \emph{morphisme} d’un graphe vers un autre est alors une fonction $f$ définie sur leurs sommets qui préserve les arêtes, au sens où, si $\tup{u,v}$ est une arête, alors $\tup{f(u),f(v)}$ doit également en être une :
nous représentons un exemple de morphisme de graphes en \Cref{fig:example-graph-homomorphism-fr}.
La notion plus riche de $\sigma$-structure, non explicitée dans ce résumé, généralise
la notion de graphe en permettant d'avoir plusieurs types de relations.

\decisionproblem{Le \textsc{problème de morphisme} sur $\sigma$}{
	Deux $\sigma$-structures $\?A$ et $\?B$.
}{
	Existe-t-il un morphisme de $\?A$ vers $\?B$ ?
}
Au-delà d'un simple problème de décision, qui se trouve être "NP"-complet,
le \textsc{problème de morphisme} sur $\sigma$ peut en fait être vu comme un 
cadre dans lequel reformuler des problèmes naturels.

\paragraph*{Bases de données.}
Considérons une requête SQL, disons la requête $\gamma$ suivante,
qui retourne les titres de films ainsi que leur horaire de projection.
\lstinputlisting[language=mysql]{fig/intro/cinema.sql}
Une base de données $\?D$ peut-être décrite comme une $\sigma$-structure---où
la signature $\sigma$ décrit le \emph{schéma} de la base---, notée également $\?D$.
Il est possible de construire%
\footnote{Voir \Cref{ex:sql-as-hom} pour plus de détails.}
de manière assez triviale---notamment en temps linéaire---une $\sigma$-structure
$\?Q_{\gamma}$, contenant notamment deux sommets \texttt{title} et \texttt{time},
tels que pour toute paire d'élements $\tup{x, y}$ de la base de données,
il existe un morphisme de $\?Q_{\gamma}$ vers $\?D$ envoyant \texttt{title}
vers $x$ et \texttt{time} vers $y$ si, et seulement si,
la paire $\tup{x,y}$ est une sortie valide de la requête $\gamma$
lorsqu'elle est évaluée sur $\?D$, c'est-à-dire que le film $x$ est projeté à
l'horaire $y$. Autrement dit, ce problème de requêtage peut-être formulé comme
un \textsc{problème de morphisme}.%
\footnote{\emph{Stricto sensu} le problème de morphisme mentionné dans ce paragraphe
ne rentre par dans le cadre décrit précédemment, puisqu'il doit gérer par constantes
(\texttt{title} et \texttt{time}). En fait, il est possible de gérer ces constantes
en les faisant passer dans la signature, et donc ce problème s'incrit en fait
bien dans le cadre étudié.}

\paragraph*{Sudoku.}
Donnons maintenant un second exemple, dans un domaine tout à fait différent.
Nous considérons une grille de Sudoku $G$, telle que celle \Cref{fig:intro-sudoku-fr}.
Cette grille peut-être encodée comme une $\sigma$-structure $\?S_{G}$,
ayant un sommet par case de la grille.%
\footnote{Voir \Cref{ex:sudoku-as-hom} pour les détails de cet encodage.}
On peut alors définir une $\sigma$-structure $\?T$ dont les sommets
sont les entiers $1, 2, \dotsc, 9$, de sorte que
les morphismes $f$ de $\?S_{G}$ vers $\?T$ décrivent exactement les complétions de la grille
de Sudoku. Notamment, étant donné un tel morphisme $f$ et une position $\tup{i,j}$,
la valeur $f(i,j)$ est la valeur rentrée dans la grille en position $\tup{i,j}$ lorsque
complétée selon $f$.
\begin{marginfigure}
	\centering%
	\begin{tikzpicture}[scale=1.3]
		\draw[step=.3333] (0, 0) grid (3, 3);
		\draw[very thick, step=1] (0, 0) grid (3, 3);

		\setcounter{row}{1}
		\setrow { }{2}{ }  {5}{ }{1}  { }{9}{ }
		\setrow {8}{ }{ }  {2}{ }{3}  { }{ }{6}
		\setrow { }{3}{ }  { }{6}{ }  { }{7}{ }

		\setrow { }{ }{1}  { }{ }{ }  {6}{ }{ }
		\setrow {5}{4}{ }  { }{ }{ }  { }{1}{9}
		\setrow { }{ }{2}  { }{ }{ }  {7}{ }{ }

		\setrow { }{9}{ }  { }{3}{ }  { }{8}{ }
		\setrow {2}{ }{ }  {8}{ }{4}  { }{ }{7}
		\setrow { }{1}{ }  {9}{ }{7}  { }{6}{ }
	\end{tikzpicture}
	\caption{
		\AP\label{fig:intro-sudoku-fr}
		Une grille de Sudoku préremplie.
		(Reproduction de \Cref{fig:intro-sudoku}.)
	}
\end{marginfigure}

Ces deux exemples illustrent que le \textsc{problème de morphisme} est versatile et peut encoder
des problèmes naturels d'une grande diversité.
Notons cependant une distinction assez notable entre les deux exemples décrits :
\begin{itemize}
	\item dans le cas des bases de données, les données sont encodées à droite du morphisme,
	et la logique à gauche : on demande s'il y a un morphisme de $\?Q_{\gamma}$ (la structure
	encodant la requête) vers $\?D$ (la base de données);
	\item au contraire, pour le Sudoku, on encode les données à gauche du morphisme, et
	la logique à droite : on demande s'il y a un morphisme de $\?S_{G}$ (la structure représentant
	la grille) vers $\?T$ (la structure représentant les contraintes du Sudoku).
\end{itemize}
Cette asymétrie entre ces deux encodages a pour conséquence que l'on, bien que ces deux exemples
partagent le même cadre, nous allons généralement étudier des questions différentes sur le
\textsc{problème de morphisme} selon qu'on soit dans le premier cas ou le second.
Il est généralement pertinent d'étudier ce qu'il se passe lorsque l'on fixe
la logique et que les données sont variables : dans le cas des bases de données,
cela revient à fixer le côté gauche du morphisme, mais dans le cas du Sudoku, cela
correspond à fixer le côté droit. Il en résulte, notamment des propriétés de monotonie et de 
complexité assez différentes. Nous séparerons le reste de cette thèse, et de ce résumé,
en deux parties.
 
\section*{Requêtes SQL \& bases de données}

L'exemple donné précédemment se généralise aisément : on peut alors montrer
que toute requête SQL de la forme
\[
	\sqlop{SELECT} \dotsc \sqlop{FROM} \dotsc \sqlop{WHERE}
\]
satisfait la même propriété, c'est-à-dire qu'évaluer une requête $\gamma$ de cette forme
sur une base de données $\?D$ revient à résoudre un \textsc{problème de morphisme} où $\gamma$ est encodé à gauche et $\?D$ à droite.
L'algorithme naïf pour résoudre le \textsc{problème de morphisme} consiste essentiellement
à énumérer chacune des
\[|\text{structure de droite}|^{|\text{structure de gauche}|}\]
fonctions allant de la structure de gauche à celle de droite, et à vérifier pour chacune si c'est 
un morphisme. Dans ce cas, la structure de droite étant une base de données, 
sachant que celles-ci peuvent faire plusieurs téraoctets en pratique, une requête de taille moyenne
peut faire exploser la complexité de cet algorithme et le rendre inefficaces.

\paragraph*{Optimisation de requêtes conjonctives.}
Une question centrale de la théorie des bases de données consiste donc à optimiser la taille 
des requêtes : c'est-à-dire, étant donné une requête, peut-on en trouver une qui est équivalente---"ie" qui retourne les mêmes réponses sur toute base de données---mais qui est plus succincte ?
Ce problème, \emph{a priori} difficile, est particulièrement bien compris dans pour les requêtes du
fragment \sqlop{SELECT}-\sqlop{FROM}-\sqlop{WHERE} de SQL, qui correspond de manière
plus abstraite au fragment existentiel positif de la logique du premier ordre,
ou encore aux \emph{requêtes conjonctives}.
Cette optimisation est précisément permise par le prisme offert par le \textsc{problème de morphisme}, et à l'encodage d'une telle requête comme une $\sigma$-structure.
Par cette méthode, nous pouvons, à partir de chaque requête conjonctive $\gamma$,
construire une requête $\widehat{\gamma}$, qui lui est sémantiquement équivalent---donc qui 
fournira les mêmes réponses sur toute base de données---et qui par ailleurs est plus petite.
De manière remarquable, ici, « plus petite » peut-être compris
au sens de « utilisant moins de variables », mais aussi dans un sens bien plus générique :
toute mesure de taille raisonnable,%
\footnote{Formellement, par raisonnable on entend « clos par sous-structure »}
est optimisée par cette requête $\widehat{\gamma}$.

\begin{figure}[h]
    \centering
    \begin{tikzpicture}
        \input{tikz/prelim-databases/ex-graph-db.tex}
    \end{tikzpicture}
    \caption{%
        \AP\label{fig:example-graph-database-fr}%
        Une base de données graphe avec deux types d'arêtes.
		(Reproduction de \Cref{fig:example-graph-database}.)
    }
\end{figure}
\paragraph*{Au-delà des requêtes conjonctives.}
Malheureusement, bien que les requêtes conjonctives soient particulièrement faciles
à optimiser, leur expressivité est relativement limitée.
Ceci est particulièrement le cas pour lorsqu'il s'agit de
requêter des bases de données graphes,%
\footnote{Appelées aussi « bases de données orientées graphe ».}
qui sont particulièrement utiles
pour stocker, par exemple, des données bibliographiques, biologiques, ou de classification.%
Une telle base de données est décrite en \Cref{fig:example-graph-database-fr}. Elle contient
deux types d'arêtes : $x \atom{\wrote} y$ indique qu'une personne $x$ a écrit un article $y$,
alors que $x \atom{\advised} y$ indique qu'une personne $x$ a supervisé la thèse de $y$.
On peut alors montrer qu'il est impossible d'exprimer la requête « l'ensemble des paires $y, z$
telles que $z$ est l'ancêtre scientifique de $y$ » avec une requête conjonctive,
ou de manière équivalente avec une requête du fragment
\sqlop{SELECT}-\sqlop{FROM}-\sqlop{WHERE} de SQL.
Pour pallier ce manque, les \emph{requêtes conjonctives à chemins réguliers}---
« conjonctive regular path queries », ou « CRPQ », en anglais---proposent de rajouter
un nouveau mécanisme dans le langage de requête :
permettre d'exiger l'existence d'un chemin allant d'un point $x$ et $y$,
et dont la suite d'étiquettes appartient à un langage régulier.
Par exemple, la CRPQ
\[
    \gamma(x, y) \defeq x \atom{\wrote} z
        \land z' \atom{\wrote} z 
        \land y \atom{({\advised})^*} z',
\]
demande toutes les paires $\tup{x,y}$ vérifiant les contraintes suivantes :
\begin{itemize}
	\item $x$ a écrit $z$,
	\item $z'$ a écrit $z$, et
	\item il y a un chemin de $y$ à $z'$ étiquetté par un mot du language
		$({\advised})^*$.
\end{itemize}
Autrement dit, cette requête demande toutes les paires de personnes $\tup{x,y}$
telles que $x$ a co-écrit un article avec un descendant scientifique de $y$.

Les CRPQ modélisent le principal mécanisme de requêtage utilisé par les langages
pour sur les bases de données graphes, comme SPARQL. L'optimisation de ces requêtes
est donc un problème primordial en théorie des bases de données.
Alors que les CRPQs étendent les requêtes conjonctives, les bonnes propriétés d'optimisation
de ces dernières ne se généralisent malheureusement pas à la classe des CRPQs, et leur
optimisation est loin d'être triviale : c'est le cœur du problème que nous étudions dans
la première moitié de cette thèse.

\begin{contribution}
	Étant donné une requête conjonctive à chemins réguliers et un entier $k$, nous pouvons décider
	s'il existe une requête équivalente utilisant au plus $k$ contraintes/atomes, et auquel cas produire
	ladite requête.
\end{contribution}

Le nombre de contraintes, ou \emph{atomes}, c'est-à-dire le nombre de termes de la
forme $x \atom{L} y$ où $L$ est un langage régulier, est une mesure pertinente de la complexité
d'évaluation d'une CRPQ. 
Cependant, nous avons également montré que, surprenamment, il existait des requêtes minimales dans 
le sens précédent---c'est-à-dire qui ne sont pas équivalentes à une CRPQ avec strictement moins 
d'atomes---, mais qui sont équivalentes à une \emph{union finie} de CRPQs utilisant chacune 
strictement moins d'atomes. Par ailleurs, l'évaluation d'une telle union étant facilement 
parallélisable, sa complexité dépend avant tout du nombre maximal d'atomes, quand bien même le 
nombre de CRPQ dans l'union serait élevé.

\begin{contribution}
	Étant donné une requête conjonctive à chemins réguliers et un entier $k$, nous pouvons décider
	s'il existe une \emph{union finie} de CRPQs qui lui est équivalente, et dont chaque
	CRPQ utilise au plus $k$ atomes. Par ailleurs, cette procédure est effective : si la réponse est oui, nous pouvons produire cette union.
\end{contribution}

Lorsque la réponse à ces problèmes est « oui », il est aisé de comprendre pourquoi : l'algorithme
produit un témoin de ce fait, c'est-à-dire une requête équivalente et plus petite.
Cependant, lorsqu'il répond par la négative, il est moins aisé de comprendre pourquoi.
Par ailleurs, mathématiquement, même sur des petites requêtes concrètes, il est difficile de
justifier pourquoi elles sont optimales.

\begin{contribution}
	Nous donnons une condition suffisante pour qu'une requête soit difficile à exprimer, 
	c'est-à-dire ne soit pas équivalente à une autre requête plus petite.
\end{contribution}

Le tour de force du résultat précédent se situe notamment dans le fait que « plus petit »
peut y être interprété de façon assez large, et ne se limite pas à la mesure du nombre d'atomes.
Il se révèle en outre particulièrement utile pour démontrer des bornes inférieures de complexité.
Malheureusement, cette condition est seulement suffisante, mais n'est pas nécessaire.
Trouver une condition qui soit nécessaire et suffisante à la minimalité d'une CRPQ est
en fait le principal problème laissé ouvert par cette thèse et diverses
pistes sont discutées en conclusion. Une telle condition, s'il s'avérait effective et efficace,
pourrait permettre de résoudre de nombreux problèmes, ou d'améliorer certains algorithmes.

Dans un second temps, nous nous intéressons à la question de la \emph{largeur arborescente}.
Il se trouve que c'est une mesure de la complexité d'une CRPQ, qui est en un sens plus fine que
le nombre d'atomes. Barceló, Romero et Vardi ont trouvé en 2016 un algorithme permettant,
étant donnée une CPRQ, de décider s'il était équivalente à une union finie
de CRPQs en forme d'arbre, c'est-à-dire de largeur arborescente $1$.
De leur propre aveu, le cas général de la largeur arborescente $k$ (pour $k$ un entier arbitraire) 
leur échappait encore.

\begin{contribution}
	Nous donnons un algorithme qui prend en entrée une requête conjonctive à chemins réguliers
	et un entier $k$, et qui décide s'il existe une union finie de CRPQs de
	largeur arborescente au plus $k$, et auquel cas la produit.
\end{contribution}

Malheureusement l'algorithme précédent, ainsi que ceux optimisant le nombre d'atomes,
sont relativement peu efficaces et fonctionnent en espace exponentiel, voire doublement exponentiel. En outre, dans chacun de ces cas, nous avons démontré des bornes de complexité
inférieures qui sont également en espace exponentiel : le problème ne vient pas tant du
fait que nos algorithmes soient peu efficaces, mais plutôt que, du fait de leur grande expressivité,
l'optimisation de CRPQ est un problème intrinsèquement difficile.
Nous nous intéressons donc à un fragment de ce langage de requête, obtenu en n'autorisant que
certains langages réguliers dans les atomes : une étude empirique récente montre que
75\% des atomes écrits en pratique sont effectivement de cette forme.
Nous donnons dans ce cas des algorithmes bien plus efficaces.

\begin{contribution}
	Nous montrons que la minimisation du nombre d'atomes (resp. de
	la largeur arborescente) d'une CRPQ n'utilisant que des expressions régulières simples
	est un problème situé dans la hiérarchie polynomiale.
\end{contribution}

En bref, les problèmes de morphismes où la logique est encodée à gauche fournit un solide
fondement à la théorie des bases de données, et à l'optimisation des requêtes conjonctives.
Dans la première partie de cette thèse, nous avons étendu une partie de ces résultats d'optimisation
aux \emph{requêtes conjonctives avec chemins réguliers} (CRPQs).

\section*{Problèmes de satisfaction de contraintes}



\end{french}